\section{Conclusion}

\paragraph{}
Ensuring an adequate food supply in response to the growing human population is crucial  \cite{carvajal-yepesGlobalSurveillanceSystem2019}, and image-based plant phenotyping has proven to be an effective tool in developing cultivars with maximum yield and resistance to various environmental stressors \cite{anandBreedingTechniquesDispense2023}. These methods are non-destructive, more time-efficient, replicable, and affordable than traditional methods \cite{daschoudhuryLeveragingImageAnalysis2019}. However, image sensors in plant phenotyping generate large amounts of data, which necessitate robust image analysis and processing tools to handle this data effectively \cite{gomesComputationalSustainabilityComputational2009}. Large plant data throughput makes image-based plant phenotyping in automated greenhouse environments cost-prohibitive, as affordable off-the-shelf hardware and sensors often lack adequate local resources to store or process large amounts of data \cite{bauckhageDataMiningPattern2013}. Consequently, remote cloud computing solutions handle image processing and analysis tasks \cite{debaucheCloudArchitecturePlant2020}, which can further increase operational costs in greenhouse management \cite{minerviniImageAnalysisNew2015}. Image-data compression techniques, low-complexity image analysis and processing techniques, and low-cost off-the-shelf hardware are effective in reducing the cost of image-based plant phenotyping by reducing the amount of computation and data throughput required \cite{minerviniApplicationawareImageCompression2013, tovarRaspberryPiPowered2018a, minerviniImageAnalysisNew2015}. However, these solutions come with marked trade-offs in terms of accuracy, precision, depth, and breadth of plant trait analysis. Moreover, existing solutions are not generalizable to other environments and plant species \cite{arnalbarbedoDigitalImageProcessing2013, liReviewImagingTechniques2014}. In short, there is a need for low-cost, open-source, flexible, modifiable, and robust plant phenotyping image processing and analysis tools that are effective and efficient on resource-constrained hardware \cite{minerviniImageAnalysisNew2015}. Future research should develop low-cost image processing tools that are precise, accurate, and adaptable to various plant monitoring contexts or parameters. In doing so, automated image-based techniques for plant monitoring will become more accessible in low-cost greenhouse environments.